\documentclass[11pt]{article}
\usepackage[margin=1in]{geometry}
\usepackage{amsmath,amssymb,amsthm}
\newtheorem*{claim}{Conjecture 00000001031}
\newtheorem*{prop}{Proposition}
\newtheorem*{thm}{Theorem}

\title{The Multiplier Group of a Singer Difference Set is not the Norm Subgroup:\\
Disproof of Conjecture 00000001031}
\author{}
\date{}

\begin{document}
\maketitle

\begin{claim}
Let $D \subset \mathbb{F}_{q^3}$ be a Singer difference set, and let its
\emph{multiplier group} be the set of multiplicative translations
$\tau$ with $\tau D = D$. Then the multiplier group is exactly the norm
subgroup $N(\mathbb{F}_{q^3}/\mathbb{F}_q)$ of $\mathbb{F}_{q^3}^*$, of index
$q^2+q+1$; this structure completely characterizes the automorphisms of
Singer difference sets, with no other multipliers.
\end{claim}

\begin{thm}
The conjecture is \textbf{false}. It already fails for $q = 2$.
\end{thm}

\section{The counterexample}

Work in $\mathbb{F}_8 = \mathbb{F}_2[\alpha]$ with $\alpha^3 = \alpha + 1$, so
that $\mathbb{F}_8^* = \langle \alpha \rangle$ is cyclic of order $7$; write
$\alpha^k \leftrightarrow k \in \mathbb{Z}_7$. The classical Singer difference
set is the set of nonzero trace-zero elements,
\[
D \;=\; \{\, x \ne 0 : \operatorname{Tr}_{8/2}(x) = x + x^2 + x^4 = 0 \,\}
\;=\; \{1,\,2,\,4\} \;\subset\; \mathbb{Z}_7 .
\]
Indeed the six ordered differences of distinct elements of $D$ are
$6,4,1,5,3,2$, i.e.\ every nonzero residue of $\mathbb{Z}_7$ occurs exactly
once, so $D$ is a $(7,3,1)$ difference set with $\lambda = q-1$.

\begin{prop}[Multiplier group]
The multipliers of $D$ in $\mathbb{Z}_7$ are exactly $\{1,2,4\}$, a group of
order $3$.
\end{prop}
\begin{proof}
Direct check for all units $\tau \in \mathbb{Z}_7^* = \{1,\dots,6\}$:
\[
\begin{array}{lll}
1\cdot D = \{1,2,4\} = D & 2\cdot D = \{2,4,1\} = D & 3\cdot D = \{3,6,5\} \ne D,\\
4\cdot D = \{4,1,2\} = D & 5\cdot D = \{5,3,6\} \ne D & 6\cdot D = \{6,5,3\} \ne D.
\end{array}
\]
Hence the multiplier group is $\{1,2,4\} = \langle 2 \bmod 7\rangle$, the
subgroup generated by the Frobenius $x \mapsto x^2$.
\end{proof}

\begin{prop}[Norm subgroup]
The norm map $N : \mathbb{F}_8^* \to \mathbb{F}_2^*$,
$N(x) = x^{1+2+4} = x^7$, has image $N(\mathbb{F}_8/\mathbb{F}_2)
= \mathbb{F}_2^* = \{1\}$: a subgroup of order $q - 1 = 1$ and index
$7 = q^2+q+1$ in $\mathbb{F}_8^*$.
\end{prop}

\begin{thm}
The multiplier group $\{1,2,4\}$ of $D$ has order $3$, while the conjectured
norm subgroup $\{1\}$ has order $q-1 = 1$. Since $\{1,2,4\} \ne \{1\}$, the
conjectured equality fails for $q = 2$. Moreover the clause ``no other
multipliers'' fails too: the multipliers $2$ and $4$ do not lie in the norm
subgroup.
\end{thm}

\section{Boundary and remarks}

Only $q = 2$ is computed; this suffices, as the conjecture is a universal
``exactly'' statement. Nothing is claimed for general $q$. The computation is
reproduced by the standalone script \texttt{reproduce.py} (polynomial
arithmetic in $\mathbb{F}_8$ from scratch) and machine-checked in core Lean 4
in the directory \texttt{lean4/} (zero axioms, zero \texttt{sorry}; run
\texttt{lake build \&\& lake env lean Check.lean}).

\end{document}
